\chapter{Albumin (Blood)}
\section*{What Is This Test?}
Serum albumin is the most abundant protein in human plasma. It makes up approximately 50--60\% of total plasma protein. It is synthesized exclusively by the liver at a rate of approximately 10--15 g/day \cite{levitt2016human}. Albumin is a protein of 66.5 kDa made of a single chain of 585 amino acids. It has three major physiologic functions. (1) It maintains plasma oncotic (colloid osmotic) pressure. Albumin contributes approximately 75--80\% of plasma oncotic pressure, which keeps fluid within the intravascular space. Hypoalbuminemia causes edema. (2) It transports numerous endogenous and exogenous ligands. These include bilirubin, fatty acids, calcium, hormones (thyroxine, cortisol), and drugs (warfarin, phenytoin, furosemide, diazepam). Approximately 45\% of total calcium is albumin-bound, so calcium must be corrected for hypoalbuminemia. (3) It has antioxidant properties (free radical scavenging at its cysteine-34 residue). Hypoalbuminemia results from four causes. First, decreased hepatic synthesis. This occurs in cirrhosis, end-stage liver disease, and malnutrition. It also occurs in acute and chronic inflammation, where the liver shifts to acute phase protein synthesis. Second, increased losses. These include nephrotic syndrome (glomerular loss), protein-losing enteropathy (GI loss), severe burns (skin loss), and hemorrhage (whole blood loss). Third, increased catabolism. This occurs in hypermetabolic states such as sepsis, burns, trauma, and malignancy. Fourth, hemodilution. This occurs in pregnancy, IV fluids, and SIADH. Hyperalbuminemia is essentially only seen in dehydration/volume contraction.
\section*{Alternative Names}
\begin{itemize}
  \item Albumin
  \item Serum Albumin
  \item ALB
  \item Plasma Albumin
\end{itemize}
\section*{Principle}
Albumin is measured by dye-binding colorimetric methods. These use bromocresol green (BCG) or bromocresol purple (BCP) on automated chemistry analyzers. BCG binds to albumin at pH 4.2 and produces a blue-green color measured spectrophotometrically. BCP is more specific for albumin with less interference from globulins and is preferred. The reference method is immunonephelometry or immunoturbidimetry.
\section*{History}
Albumin is named from the Latin albumen (egg white). It was among the first proteins studied by 19th century protein chemists. Egg albumin was first crystallized by Franz Hofmeister in 1890. In the 1940s, Edwin J. Cohn and colleagues at Harvard developed cold ethanol (Cohn) fractionation of plasma proteins \cite{peters1996all}. This isolated serum albumin as fraction V. It made albumin available as a plasma volume expander for the treatment of shock and burns. The modern clinical measurement of serum albumin began in 1971. In that year, Doumas, Watson, and Biggs described the bromocresol green dye-binding method \cite{doumas1971albumin}. This method remains the basis of automated albumin assays. Albumin has since become a marker of hepatic synthetic function and nutritional status \cite{vincent2003hypoalbuminemia}. When low, it is an independent predictor of morbidity and mortality in hospitalized patients.
\section*{Why This Test Is Ordered}
\begin{itemize}
  \item Evaluation of liver disease and hepatic synthetic function, as part of a liver panel and the Child-Pugh score.
  \item Assessment of nutritional status, especially in hospitalized, postoperative, or malnourished patients.
  \item Investigation of edema, ascites, or anasarca, and evaluation for nephrotic syndrome (with urine protein).
  \item Workup of suspected protein-losing enteropathy, burns, or chronic inflammation.
  \item Preoperative risk assessment, since hypoalbuminemia predicts postoperative complications.
  \item Correction of serum calcium, which must be adjusted for the albumin level before interpretation.
\end{itemize}
\section*{Primary vs Secondary Test}
Serum albumin is a secondary and monitoring test rather than a primary diagnostic test. It is most often ordered as part of a liver function panel or comprehensive metabolic panel. These assess hepatic synthesis and nutritional status. When abnormal, it directs further evaluation (urine protein, serum protein electrophoresis, stool alpha-1 antitrypsin) to identify the cause.
\section*{How This Test Is Performed}
A blood sample is collected by standard venipuncture into a serum separator (gold-top) or plain red-top tube. In the laboratory, albumin is measured by the bromocresol green or bromocresol purple dye-binding method. It is measured on an automated chemistry analyzer, usually as part of a comprehensive metabolic panel. Results are typically available within a few hours of collection.
\section*{What Preparation Is Needed}
No fasting or special preparation is required. The sample is often drawn with morning fasting chemistry panels. Posture affects the result. Standing for 15 minutes or longer can raise the serum albumin concentration by up to 10\%. A consistent sampling position is advisable when serial measurements are compared. Prolonged tourniquet use should be avoided.
\section*{Contraindications and Precautions}
\begin{itemize}
  \item There are no absolute contraindications to standard venipuncture. Several interpretive cautions apply. Albumin is a negative acute phase reactant. Its levels fall during acute inflammation, surgery, and critical illness regardless of nutritional state \cite{francharcas2001meaning}. Hypoalbuminemia in hospitalized patients often reflects the acute phase response and fluid shifts rather than true deficiency \cite{vincent2003hypoalbuminemia}. The bromocresol green method can overestimate albumin in the presence of paraproteins and in some patients with renal or hepatic disease. In these cases, confirmatory measurement with bromocresol purple or immunonephelometry may be required. Total serum calcium must be corrected for the albumin concentration before interpretation.
\end{itemize}
\section*{Risks}
Minimal risk. Slight pain or bruising at the venipuncture site. Rarely: hematoma, infection, or vasovagal syncope.
\section*{What Happens After the Test}
Results are usually available the same day. They are interpreted with total protein, liver enzymes, and the clinical history. A low albumin prompts investigation of the underlying cause. These causes include hepatic synthetic failure, protein losses (renal, gastrointestinal, or cutaneous), malnutrition, or inflammation. Serial measurements can monitor nutritional support or the course of disease.
\section*{Understanding Results}
\subsection*{Hypoalbuminemia (<3.5 g/dL)}
Decreased synthesis: cirrhosis, chronic liver disease, and malnutrition (protein-calorie malnutrition, kwashiorkor). Decreased synthesis also occurs in acute and chronic inflammation (negative acute phase reactant). Increased loss: nephrotic syndrome (urine protein >3.5 g/24 hours), protein-losing enteropathy, severe burns, and hemorrhage. Dilutional: pregnancy, IV fluid administration, and SIADH. Serum calcium correction: corrected calcium (mg/dL) = measured calcium + [0.8 $\times$ (4.0 - albumin)].
\subsection*{Hyperalbuminemia (>5.0 g/dL)}
Dehydration/volume contraction (hemoconcentration). No known pathologic overproduction syndrome.
\section*{Normal Results}
Serum albumin: 3.5--5.0 g/dL (35--50 g/L). Newborns: 2.8--4.4 g/dL. This is lower, and it reaches adult levels by 3--6 months. Hypoalbuminemia is an important prognostic marker. Albumin <3.5 g/dL is linked to increased mortality, length of hospitalization, surgical complications, and poor wound healing. Serum calcium must be corrected for hypoalbuminemia.
\section*{Follow-up Tests}
Serum total protein. Serum protein electrophoresis (SPEP). Liver function tests (AST, ALT, ALP, bilirubin, PT/INR). Prealbumin (transthyretin). It has a shorter half-life of 2 days versus 20 days for albumin, so it is a better marker of recent nutritional status. Urine protein (24-hour or spot UPCR). Stool alpha-1 antitrypsin clearance (protein-losing enteropathy).
\section*{References}
\bibliographystyle{unsrtnat}
\bibliography{references}

\clearpage
\section*{Regulatory Status and Authorization Date}
This chapter describes a test category, not a specific commercial product. FDA, CDSCO, EU IVDR, and MHRA approval or registration dates therefore cannot be assigned without the assay manufacturer, product name, instrument, and intended use. For laboratory-developed tests, an individual agency approval date may not exist. See the Preface for the interpretation of regulatory status.

\clearpage
